\documentclass[10pt,landscape,a4paper]{article}
\usepackage[landscape,margin=1.2cm]{geometry}
\usepackage{multicol}
\usepackage{tabularx}
\usepackage{xcolor}
\usepackage{hyperref}
\usepackage[T1]{fontenc}
\usepackage{lmodern}
\usepackage{amssymb}
\usepackage{fuzz}
\usepackage{proof}

% Compact spacing
\setlength{\parindent}{0pt}
\setlength{\parskip}{2pt plus 1pt}
\setlength{\columnsep}{1.2cm}
\pagestyle{empty}

% Section formatting
\usepackage{titlesec}
\titleformat{\section}{\large\bfseries\color{darkgray}}{}{0pt}{}[\vspace{-4pt}\rule{\columnwidth}{0.4pt}]
\titlespacing{\section}{0pt}{6pt}{2pt}

% Code style
\newcommand{\kw}[1]{\texttt{\bfseries #1}}
\newcommand{\op}[1]{\texttt{#1}}
\newcommand{\renders}{\,$\rightarrow$\,}

% Stacked example: input then rendered output (compact)
\newenvironment{exwrite}%
  {\par\noindent\scriptsize\textbf{You write:}\par\vspace{0pt}\scriptsize}%
  {\vspace{1pt}}
\newenvironment{exget}%
  {\par\noindent\scriptsize\textbf{You get:}\par\vspace{0pt}\small}%
  {\par\vspace{3pt}}

% Table helper — left col monospace, right col roman
\newcolumntype{L}{>{\ttfamily}l}
\newcolumntype{R}{>{\raggedright\arraybackslash}X}

\begin{document}

\begin{center}
{\LARGE\bfseries txt2tex Reference}\\[2pt]
{\small Write math like you're at a whiteboard.  Get \LaTeX\ you can hand in.}\\[1pt]
{\footnotesize\color{gray}\url{https://github.com/jmf-pobox/txt2tex}}
\end{center}

\vspace{-2pt}

\begin{multicols}{3}
\raggedcolumns

% ── DOCUMENT STRUCTURE ──────────────────────────
\section{Document Structure}

\begin{tabularx}{\columnwidth}{LR}
=== Title ===        & \verb|\section*{Title}| \\
** Solution N **     & Bold heading with spacing \\
(a) Text             & Part label \\
TEXT: paragraph       & Prose block (smart formula detection) \\
LATEX: \textbackslash cmd & Raw \LaTeX\ passthrough \\
PAGEBREAK:           & Page break \\
\end{tabularx}

\smallskip
\begin{tabularx}{\columnwidth}{LR}
TITLE: ...           & Document title \\
AUTHOR: ...          & Author name \\
DATE: ...            & Date string \\
BIBLIOGRAPHY: f.bib  & BibTeX file \\
{[cite key]}         & \verb|\citep{key}| \\
\end{tabularx}

% ── IDENTIFIERS ─────────────────────────────────
\section{Identifier Decoration}

\begin{tabularx}{\columnwidth}{LR}
count'               & $count'$ (after-state) \\
in?                  & $in?$ (input) \\
out!                 & $out!$ (output) \\
x?'                  & $x?'$ (runs may combine) \\
'sunk'               & $\mbox{`sunk'}$ (string literal) \\
\end{tabularx}

{\footnotesize Decoration attaches to any identifier. Reserved words (\texttt{theta}, \texttt{defs}, \texttt{hide}, \texttt{project}, \texttt{Delta}, \texttt{Xi}) cannot be decorated.}

% ── LOGIC ───────────────────────────────────────
\section{Logic}

\begin{tabularx}{\columnwidth}{LR}
p land q             & $p \land q$ \\
p lor q              & $p \lor q$ \\
lnot p               & $\lnot p$ \\
p => q               & $p \Rightarrow q$ \\
p <=> q              & $p \Leftrightarrow q$ \\
forall x : S | P     & $\forall x : S @ P$ \\
exists x : S | P     & $\exists x : S @ P$ \\
exists\_1 x : S | P  & $\exists_1 x : S @ P$ \\
lambda x : S | E     & $\lambda x : S @ E$ \\
mu x : S | P         & $\mu x : S @ P$ \\
\end{tabularx}

\smallskip
{\footnotesize\textbf{Note:} Use \kw{land}, \kw{lor}, \kw{lnot}. English \texttt{and}/\texttt{or}/\texttt{not} are not supported.}

% ── SETS ────────────────────────────────────────
\section{Sets \& Types}

\begin{tabularx}{\columnwidth}{LR}
x elem S             & $x \in S$ \\
x notin S            & $x \notin S$ \\
A subset B           & $A \subseteq B$ \\
A union B            & $A \cup B$ \\
A inter B            & $A \cap B$ \\
A setminus B         & $A \setminus B$ \\
power S              & $\power S$ \\
F S                  & $\finset S$ \\
\{x : S | P\}       & Set comprehension \\
A cross B            & $A \cross B$ \\
N / Z / N1           & $\nat$ / $\num$ / $\nat_1$ \\
n..m                 & $n \upto m$ \\
\end{tabularx}

% ── RELATIONS ───────────────────────────────────
\section{Relations}

\begin{tabularx}{\columnwidth}{LR}
A <-> B              & $A \rel B$ \\
x |-> y              & $(x \mapsto y)$ \\
dom R / ran R        & $\dom R$ / $\ran R$ \\
A dres R             & $A \dres R$ \\
A dsub R             & $A \ndres R$ \\
R rres B             & $R \rres B$ \\
R rsub B             & $R \nrres B$ \\
R oplus S            & $R \oplus S$ (override) \\
R o9 S               & $R \semi S$ (forward comp.) \\
R comp S             & $R \comp S$ (backward comp.) \\
\end{tabularx}

% ── FUNCTIONS ───────────────────────────────────
\section{Functions}

\begin{tabularx}{\columnwidth}{LR}
A pfun B             & $A \pfun B$ \\
A fun B              & $A \fun B$ (total) \\
A pinj B             & $A \pinj B$ \\
A inj B              & $A \inj B$ \\
A surj B             & $A \surj B$ \\
A bij B              & $A \bij B$ \\
A ffun B             & $A \ffun B$ (finite) \\
f(x)                 & Function application \\
\end{tabularx}

% ── SEQUENCES ───────────────────────────────────
\section{Sequences}

\begin{tabularx}{\columnwidth}{LR}
<> / <a, b>          & $\langle\rangle$ / $\langle a, b \rangle$ \\
s \textasciicircum{} t  & $s \cat t$ (concatenation) \\
seq S / seq1 S       & $\seq S$ / $\seq_1 S$ \\
\# s                 & $\# s$ (length) \\
\end{tabularx}

% ── Z NOTATION BLOCKS ───────────────────────────
\section{Z Notation Blocks}

\begin{tabularx}{\columnwidth}{LR}
given A, B           & Given types: $[A, B]$ \\
T ::= a | b          & Free type \\
name == expr         & Abbreviation \\
\end{tabularx}

\smallskip
\begin{tabularx}{\columnwidth}{LR}
\multicolumn{2}{l}{\texttt{axdef} / \texttt{schema Name} / \texttt{gendef [X]}}\\
\quad declarations   & Variable declarations \\
\texttt{where}       & Separator \\
\quad predicates     & Constraining predicates \\
\texttt{end}         & Close the block \\
\end{tabularx}

% ── SCHEMA INCLUSION & CALCULUS ─────────────────
\section{Schema Inclusion \& Calculus}

\begin{tabularx}{\columnwidth}{LR}
Delta Counter        & $\Delta Counter$ (before/after) \\
Xi Card              & $\Xi Card$ (read-only) \\
SName                & bare inclusion \\
theta S              & $\theta S$ (state binding) \\
theta S'             & $\theta S'$ (after-state) \\
\end{tabularx}

\smallskip
\begin{tabularx}{\columnwidth}{LR}
\multicolumn{2}{l}{\texttt{Name defs RHS} — horizontal definition:}\\
Name defs Schema     & $Name \defs Schema$ \\
N defs Delta C       & $N \defs \Delta C$ \\
N defs [d : T | P]  & inline schema text \\
N[X] defs G[X]      & generic form \\
\end{tabularx}

\smallskip
\begin{tabularx}{\columnwidth}{LR}
\multicolumn{2}{l}{\texttt{defs} RHS operators (schema calculus):}\\
S[old/new]           & rename component \\
S[a/b, c/d]         & multiple renames \\
S ; T                & $S \semi T$ (composition) \\
S >> T               & $S \pipe T$ (piping) \\
S hide (x, y)        & $S \hide (x,y)$ \\
S project T          & $S \project T$ \\
\end{tabularx}

% ── PROOF SYNTAX REFERENCE ──────────────────────
\section{Proof Tree Syntax}

\begin{tabularx}{\columnwidth}{LR}
\textrm{indent}      & Premises indented under conclusion \\
{[rule]}              & Justification label \\
{[rule {[n]}]}        & Discharge assumption $n$ \\
{[n] expr}            & Boxed assumption labelled $n$ \\
::                    & Sibling premises \\
\end{tabularx}

\smallskip
{\footnotesize Rules: \texttt{land intro}, \texttt{land elim 1/2}, \texttt{lor intro 1/2}, \texttt{lor elim}, \texttt{=> intro/elim}, \texttt{<=> intro/elim}, \texttt{lnot intro/elim}, \texttt{forall intro/elim}, \texttt{exists intro/elim}.}

% ── USAGE ───────────────────────────────────────
\section{Usage}

\begin{tabularx}{\columnwidth}{LR}
txt2tex f.txt        & $\rightarrow$ f.tex + f.pdf \\
txt2tex f.txt {-}{-}tex-only & $\rightarrow$ f.tex only \\
txt2tex -i           & Interactive REPL \\
txt2tex {-}{-}check-env  & Verify dependencies \\
\end{tabularx}

\end{multicols}

%% ── PAGE 2: PROOFS BY EXAMPLE ──────────────────
\newpage

\begin{center}
{\LARGE\bfseries txt2tex Proofs --- By Example}\\[2pt]
{\small Each example shows what you type and what txt2tex renders.}
\end{center}

\vspace{-2pt}

\begin{multicols}{3}
\raggedcolumns

% ── HOW PROOFS WORK ─────────────────────────────
\section{How Proof Trees Work}

{\scriptsize
In txt2tex, proof trees are written \textbf{bottom-up}: the conclusion is the first line, and its premises are indented below it. Think of it as reading the inference rule from conclusion to justification:

\smallskip
\begin{tabularx}{\columnwidth}{LR}
\textrm{conclusion}  & First line (least indented) \\
\textrm{\quad premise}& Indented = supports the line above \\
{[rule]}              & Justification in brackets \\
{[n] expr}            & Boxed assumption (labelled $n$) \\
{[rule {[n]}]}        & Discharge assumption $n$ \\
::                    & Marks sibling premises \\
\end{tabularx}

\smallskip
Deeper indentation = deeper in the tree. Two premises at the same indent level are siblings that jointly support their parent.
}

% ══ PROPOSITIONAL LOGIC (Lecture 1) ═════════════

% ── TRUTH TABLE ─────────────────────────────────
\section{Truth Table {\normalsize\color{gray}\textrm{\scriptsize Prop.\ Logic}}}

\begin{exwrite}%
\begin{verbatim}
TRUTH TABLE:
p | q | p => q
T | T | T
T | F | F
F | T | T
F | F | T
\end{verbatim}
\end{exwrite}
\begin{exget}%
\begin{center}
\begin{tabular}{ccc}
$p$ & $q$ & $p \Rightarrow q$ \\ \hline
T & T & T \\
T & F & F \\
F & T & T \\
F & F & T \\
\end{tabular}
\end{center}
\end{exget}

% ── EQUIV CHAIN (De Morgan) ────────────────────
\section{Equivalence Chain {\normalsize\color{gray}\textrm{\scriptsize Prop.\ Logic}}}

\begin{exwrite}%
\begin{verbatim}
EQUIV:
lnot (p land q)
lnot p lor lnot q  [De Morgan]
\end{verbatim}
\end{exwrite}
\begin{exget}%
\begin{center}
$\begin{array}{l@{\hspace{2em}}r}
\lnot (p \land q) \\
\Leftrightarrow \lnot p \lor \lnot q
  & [\mbox{De Morgan}]
\end{array}$
\end{center}
\end{exget}

{\scriptsize \kw{EQUIV:} joins steps with $\Leftrightarrow$. Use for \emph{predicate} equivalences.}

% ══ PROOF TREES (Lecture 4) ═════════════════════

% ── MODUS PONENS ────────────────────────────────
\section{Modus Ponens {\normalsize\color{gray}\textrm{\scriptsize Proof Trees}}}

\begin{exwrite}%
\begin{verbatim}
PROOF:
  q [=> elim]
    p => q
    p
\end{verbatim}
\end{exwrite}
\begin{exget}%
\begin{center}
$\infer[\Rightarrow\textrm{-elim}]{q}{p \Rightarrow q & p}$
\end{center}
\end{exget}

% ── AND-INTRO ───────────────────────────────────
\section{$\land$-intro {\normalsize\color{gray}\textrm{\scriptsize Proof Trees}}}

\begin{exwrite}%
\begin{verbatim}
PROOF:
  p land q [land intro]
    p
    q
\end{verbatim}
\end{exwrite}
\begin{exget}%
\begin{center}
$\infer[\land\textrm{-intro}]{p \land q}{p & q}$
\end{center}
\end{exget}

% ── IMPLIES-INTRO WITH DISCHARGE ────────────────
\section{$\Rightarrow$-intro with discharge {\normalsize\color{gray}\textrm{\scriptsize Proof Trees}}}

\begin{exwrite}%
\begin{verbatim}
PROOF:
  p => p [=> intro [1]]
    [1] p
\end{verbatim}
\end{exwrite}
\begin{exget}%
\begin{center}
$\infer[\Rightarrow\textrm{-intro}^{[1]}]{p \Rightarrow p}{\ulcorner p \urcorner^{[1]}}$
\end{center}
\end{exget}

{\scriptsize \texttt{[1]~p} = boxed assumption. \texttt{[=> intro~[1]]} discharges it.}

% ── LOR-ELIM (CASE ANALYSIS) ───────────────────
\section{$\lor$-elim / case analysis {\normalsize\color{gray}\textrm{\scriptsize Proof Trees}}}

\begin{exwrite}%
\begin{verbatim}
PROOF:
  r [lor elim]
    p lor q
    :: p => r
      [1] p
        r [from p]
    :: q => r
      [2] q
        r [from q]
\end{verbatim}
\end{exwrite}
\begin{exget}%
\begin{center}
$\infer[\lor\textrm{-elim}]{r}{p \lor q & \infer{r}{\ulcorner p \urcorner^{[1]}} & \infer{r}{\ulcorner q \urcorner^{[2]}}}$
\end{center}
\end{exget}

{\scriptsize \texttt{::} = sibling premises (case arms).}

% ══ INDUCTION / EQUATIONAL (Lectures 11–13) ════

% ── EQUAL CHAIN ─────────────────────────────────
\section{Equality Chain {\normalsize\color{gray}\textrm{\scriptsize Induction}}}

\begin{exwrite}%
\begin{verbatim}
EQUAL:
0 + 0
0      [0+0 = 0]
length(<>)
       [length(<>) = 0]
\end{verbatim}
\end{exwrite}
\begin{exget}%
\begin{center}
$\begin{array}{l@{\hspace{1.5em}}r}
0 + 0 \\
= 0 & [\mbox{0+0 = 0}] \\
= length(\langle\rangle) & [\mbox{length}(\langle\rangle) = 0]
\end{array}$
\end{center}
\end{exget}

{\scriptsize \kw{EQUAL:} = expression chains ($=$). Use for $\mathbb{N}$/sets/sequences.}

\end{multicols}

%% ── PAGE 3: RELATIONAL DATABASES ──────────────
\newpage

\begin{center}
{\LARGE\bfseries txt2tex Relational Databases --- By Example}\\[2pt]
{\small DAT course notation: primary keys, relational algebra, bindings, GROUP/UNGROUP.}
\end{center}

\vspace{-2pt}

\begin{multicols}{3}
\raggedcolumns

% ── PRIMARY KEYS ────────────────────────────────
\section{Primary Keys}

{\scriptsize Mark a primary-key attribute with \kw{pk} in a named schema body.
The PK annotation sits \emph{outside} the Z box so fuzz can type-check the schema cleanly.}

\begin{exwrite}
\begin{verbatim}
schema Class
  pk class : ClassName
  country  : CountryName
  bore     : N
end
\end{verbatim}
\end{exwrite}
\begin{exget}
\begin{center}
\begin{schema}{Class}
class : ClassName \\
country : CountryName \\
bore : \nat
\end{schema}\\[2pt]
\noindent$\mathrm{PK}(\mathrm{Class}) = \{class\}$
\end{center}
\end{exget}

{\scriptsize Composite: two \kw{pk} lines $\Rightarrow$ $\mathrm{PK}(S)=\{a,b\}$.
Comma-separated names on one \kw{pk} line also work.
\kw{pk} is rejected in \texttt{axdef} / \texttt{gendef} / inline schema text.}

% ── RELATIONAL ALGEBRA ──────────────────────────
\section{Relational Algebra}

\begin{tabularx}{\columnwidth}{LR}
sigma[p](R)          & $\mathrm{Restrict}_{p}(R)$ — restrict \\
pi[A,B](R)           & $\mathrm{Project}\{A,B\}(R)$ — project \\
rho[A as B](R)       & $\mathrm{Rename}_{A \to B}(R)$ — rename \\
R bowtie S           & $R \otimes S$ — natural join \\
R bowtie [p] S       & $\mathrm{Join}_{p}(R, S)$ — theta-join \\
R div S              & $R \div S$ — division \\
\end{tabularx}

\smallskip
{\footnotesize \texttt{sigma}/\texttt{pi}/\texttt{rho} bind at atom level (prefix + args). \texttt{bowtie} and \texttt{div} are infix at cross-product precedence.}

\smallskip
{\scriptsize\textbf{Naming a query.} Use \texttt{Name == expr} (Z abbreviation) for both pure-Z and DAT-bearing right-hand sides. txt2tex inspects the RHS: when it contains a DAT construct (algebra, binding, GROUP/UNGROUP), the abbreviation emits as \texttt{\textbackslash noindent\$\ldots\$} outside any Z block; otherwise it emits inside a \texttt{zed} block.}

\smallskip
{\scriptsize\textbf{Fuzz note:} algebra expressions emit \emph{outside} any Z environment, so fuzz silently skips them — your schemas and axdefs in the same document still type-check cleanly. Do \emph{not} wrap algebra inside \texttt{zed}/\texttt{axdef}/\texttt{schema} blocks; fuzz parses block contents strictly as Z and will reject the algebra subscripts.}

\begin{exwrite}
\begin{verbatim}
BigGuns == pi[class,country](
  sigma[bore >= 16](Class))
\end{verbatim}
\end{exwrite}
\begin{exget}
\begin{center}
\noindent$BigGuns \defs \mathrm{Project}\{class,country\}(\mathrm{Restrict}_{bore \geq 16}(Class))$
\end{center}
\end{exget}

% ── BINDINGS ────────────────────────────────────
\section{Z Bindings \& Set Comprehension}

{\scriptsize Z RM §3.7 binding brackets construct labelled tuples.}

\begin{tabularx}{\columnwidth}{LR}
\{| name == e |\}    & $\lblot name == e \rblot$ \\
\{| a == e, b == f |\} & multiple components \\
\{| |\}              & empty binding \\
\end{tabularx}

\smallskip
{\scriptsize Multi-typed comprehension with binding:}

\begin{exwrite}
\begin{verbatim}
{ s : Ship; c : Class |
  s.class = c.class .
  {| name == s.name |} }
\end{verbatim}
\end{exwrite}
\begin{exget}
\begin{center}
$\{ s : Ship; c : Class \mid s.class = c.class \bullet \lblot name == s.name \rblot \}$
\end{center}
\end{exget}

{\scriptsize Use \texttt{;} to separate variable-type pairs inside \texttt{\{ \}}. Binding components are \textbf{comma}-separated (Z RM §3.7).}

\smallskip
{\scriptsize\textbf{Fuzz note:} bindings emit \emph{outside} any Z environment, so fuzz silently skips them. Schemas and axdefs in the same document still type-check. Do \emph{not} place a binding inside a \texttt{zed}/\texttt{axdef}/\texttt{schema} block; fuzz parses block contents strictly and rejects \texttt{==} inside $\lblot \ldots \rblot$ even though both the brackets and the operator are defined in \texttt{fuzz.sty}.}

% ── GROUP / UNGROUP ──────────────────────────────
\section{GROUP \& UNGROUP}

{\scriptsize Date's nested-relation operators.}

\begin{tabularx}{\columnwidth}{LR}
R group (\{A,B\} as alias) & $R\ \mathop{\mathrm{GROUP}} (\{A,B\}\ \mathop{\mathrm{AS}}\ alias)$ \\
R ungroup alias    & $R\ \mathop{\mathrm{UNGROUP}}\ alias$ \\
\end{tabularx}

\begin{exwrite}
\begin{verbatim}
Members group ({name,email} as contact)
Members ungroup contact
\end{verbatim}
\end{exwrite}
\begin{exget}
\begin{center}
$Members\ \mathop{\mathrm{GROUP}}\ (\{name,email\}\ \mathop{\mathrm{AS}}\ contact)$\\[3pt]
$Members\ \mathop{\mathrm{UNGROUP}}\ contact$
\end{center}
\end{exget}

{\scriptsize\textbf{Fuzz note:} GROUP/UNGROUP emit \emph{outside} any Z environment, so fuzz silently skips them. Do \emph{not} place these inside \texttt{zed}/\texttt{axdef}/\texttt{schema} blocks — the \texttt{\textbackslash mathop} wrapping and Date braces sit outside Z's grammar and fuzz rejects them in block contents.}

\begin{exget}
\begin{center}
\end{center}
\end{exget}

{\scriptsize \texttt{group}/\texttt{ungroup} use \texttt{\textbackslash mathop\{\textbackslash mathrm\{...\}\}} — kernel LaTeX, no extra package.}

% ── QUICK REFERENCE ─────────────────────────────
\section{DAT Quick Reference}

\begin{tabularx}{\columnwidth}{LR}
pk a : T             & primary key in schema \\
pk a, b : T          & composite pk \\
sigma[p](R)          & restrict \\
pi[A](R)             & project \\
rho[A as B](R)       & rename attr \\
R bowtie S           & natural join ($\otimes$) \\
R div S              & division \\
\{| a == e |\}       & binding bracket \\
\{S; C | P . E\}    & multi-typed set comp \\
R group (\{A\} as x) & group attrs \\
R ungroup x          & ungroup attr \\
\end{tabularx}

\end{multicols}

\end{document}
